\section{Methodology}
\label{sec:method}

The evaluation has two vehicles. A replay harness compares the ledger with the two passive detectors on one shared stream of sprayed traffic, because the passive detectors are host-side systems that cannot run on our fabric and because a head-to-head needs every arm to see the same packets. A silicon testbed then measures the ledger's own side of that comparison on a real switch. In this section, we describe the harness, the three arms, the metrics, and the testbed.

\subsection{The Replay Harness}

\textbf{Shared stream.} Every epoch, the harness draws one per-packet spray decision and one per-hop survival outcome for every packet, and it feeds the same draw to all three arms. Spray is independent and uniform across the $k=8$ spines, which matches SprayCheck's own model and FlowPulse's temporal-symmetry assumption~\cite{spraycheck,flowpulse}. Survival is independent per directed link with rate $1-p$ on a faulty link and $1-10^{-5}$ on a healthy one. Each arm sees only what its own switch would see. The ledger arm receives the exact per-sublink transmit and receive counts its registers produce. The SprayCheck arm receives receive-only per-spine arrival counts; no transmit count and no drop count exists in its interface, so none can be passed to it by mistake. The FlowPulse arm receives receive-only per-port load, predicted from its own bootstrapped history, and never the ledger's transmit ground truth.

\textbf{Scenario.} Each trial runs ten epochs of clean traffic, during which FlowPulse bootstraps its baseline and the ledger warms its floor. The trial then degrades one spine, or one directed link, to the target rate $p$ for the rest of the run. For detection (O1) the fleet carries two million packets per epoch, and a trial ends when the arm flags the faulty spine or after 80 post-onset epochs, a budget of 160 million packets. For localization (O2) the harness extends to the two-hop fabric of $n=4$ leaves, and each ordered leaf pair carries two million packets per epoch. The localization budget is 60 post-onset epochs, which is deliberately generous to the passive arms so that SprayCheck can complete a corroborating round of measured flows. For the correlated gate (O4) the window is 40 post-onset epochs. We run 50 seeds per cell.

\textbf{Two fidelity checks that changed the harness.} Before trusting any number we checked each arm against its own paper. Two checks failed and were fixed. First, a per-packet traffic generator made one two-million-packet epoch take tens of seconds; we vectorized it with a multinomial draw, which changed no distribution. Second, at an initial volume of 200 thousand packets per epoch, FlowPulse's fixed one-percent threshold false-alarmed on a healthy spine in every trial. The reason is that the relative spray noise at that volume, about 0.6 percent, sat close to its threshold. We raised the volume to two million packets per epoch, at which its false-positive rate on healthy spines is 0.0000, consistent with the paper's own claim of under one percent~\cite{flowpulse}. We report that rate next to every result.

\subsection{The Three Arms}

\textbf{The ledger.} The controller of Section~\ref{sec:ledger}, unmodified from the version measured on silicon, runs over the exact per-sublink counts. Its localized set is the set of sublinks e-BH rejects.

\textbf{SprayCheck-Z.} We reimplement SprayCheck's detector from its paper~\cite{spraycheck}: a Z-test on per-spine arrival counts with sensitivity $s$ calibrated under the paper's own Poisson model at 2.5 million packets per spine, which gives $s=3.24$. For localization we implement its full cross-leaf rule: a report flags the uplink and downlink pair of a path, and a single link is named only when it lies in the intersection of reports from different leaves. SprayCheck's real spraying uses join-the-shortest-of-two-queues, which has less noise than the Poisson model; a sensitivity tuned to that would shift its cost curve down by a constant factor but not change how the cost scales with $p$.

\textbf{FlowPulse-$\theta$.} We reimplement FlowPulse's port-load detector with its fixed one-percent threshold against a learned per-port baseline, and its per-sender localization rule: if every sender on an ingress port is affected the local downlink is named, and if only one sender is affected the uplink from that sender's leaf is named~\cite{flowpulse}. We give the arm per-sender visibility on every port, which its paper assumes; this is the charitable reading, not a weakened one.

\subsection{Metrics}

\textbf{Action rate} is the fraction of the 50 seeds in which the arm flags the faulty spine or link within the budget, with a 95 percent Wilson interval~\cite{wilson}. \textbf{Packets to detect} is the fleet-wide packet count from the fault's onset to the first correct flag, reported as a median with its interquartile range over the seeds that detected; we mark it as survivorship-biased wherever the action rate is below one. \textbf{False positive} is any healthy spine or link flagged, reported next to every action rate. For O2, \textbf{exact} means the arm named the true directed link and only it, \textbf{ambiguous} means it named a set that contains a healthy link, and \textbf{miss} means it named nothing. We also report the mean set size over detections and a paired exact McNemar test between arms on the same seeds. For O4, \textbf{recall} is the fraction of true faulty links named, and \textbf{exact-all} is the fraction of seeds that named the true set and nothing else. We report the false-positive rate both over the whole window and at its final epoch, so that a transient false flag can be separated from a persistent one.

\subsection{The Silicon Testbed}

\textbf{Switch.} The ledger runs on an Intel Tofino~1 switch~\cite{tofino} with SDE 9.13.2. A 4-leaf, 2-spine virtual fabric is emulated on the switch itself: a four-lane 25~Gb/s direct-attach cable loops four front-panel ports back to the switch, and each directed virtual link is one traffic-manager queue on one loop port. With four traffic classes, the program tracks 1024 sublinks. The pipeline occupies 12 ingress and 5 egress match-action stages (Section~\ref{sec:cost}).

\textbf{Traffic and ground truth.} A server with a 25~Gb/s network interface sends 1400-byte packets at 20 thousand packets per second, and a differentiated-services code point selects the traffic class and thus the sublink. Two fault injectors run in the switch's egress pipeline. The exact injector drops exactly $D$ stamped packets on a chosen sublink, so the ground truth is $D$. The stochastic injector draws a 16-bit random value per packet and drops the packet when the value falls below $W=\lfloor p\cdot 2^{16}\rceil$, so it drops each packet with probability $W/2^{16}$. For this injector the ground truth is its own direct counter of packets dropped, read from the switch, because a random injector has no preset count. Every cell reads a baseline census immediately before arming and reports deltas against it. Every cell also keeps $\Delta\mathit{seq}$ below $2^{16}$, so that the sequence wraps at most once, and clears both injectors before and after. The recovered loss is Equation~\eqref{eq:loss} read through the controller's census.

\textbf{Model validation.} Before any hardware run, the compiled program passed nine correctness tests on the vendor's software model. They cover the exact recovery of injected loss, the invisibility of a trailing loss until the next survivor, the 32-bit transmit frontier past 255 and 256 packets, the absence of phantom loss under reordering, the $-1$ signature of a duplicate, and the realized rate of the stochastic injector, which fell within four standard deviations of its target. The hardware measurements of Section~\ref{sec:silicon} were then taken on the same build.
